\documentclass[10pt,a4paper]{amsart}
\usepackage[margin=19mm]{geometry}
\usepackage{amsmath,amssymb,mathtools}
\usepackage[scr=rsfs,cal=euler]{mathalfa}
\usepackage{xfrac}
\usepackage[unicode,hidelinks,bookmarks=false]{hyperref}
\newcommand{\ie}{i.e.\ }
\newcommand{\cf}{cf.\ }
\newcommand{\resp}{respectively\ }
\newcommand{\Subsection}{\subsection}
\providecommand{\nobreakdash}{\nobreak\hbox{-}}
\setlength{\emergencystretch}{2em}
\multlinegap0pt
\usepackage{bm}
\usepackage{bbm}
\usepackage{dsfont}
\usepackage{xspace}
\usepackage{cancel}
\usepackage{mathtools}
\usepackage{amsthm}
\makeatletter
\@ifundefined{theorem}{\newtheorem{theorem}{Theorem}}{}
\@ifundefined{prop}{\newtheorem{prop}[theorem]{Proposition}}{}
\makeatother
\makeatletter
\expandafter\ifx\csname cn\endcsname\relax
\newenvironment{cn}{\noindent\color{blue} }{}
\fi
\expandafter\ifx\csname edit\endcsname\relax
\newenvironment{edit}{\color{OliveGreen} }{}
\fi
\expandafter\ifx\csname ok\endcsname\relax
\newenvironment{ok}{\noindent\color{red} }{}
\fi
\makeatother
\title[Equations in Products of Free Groups and 3-Manifold Groups, ]{Equations in Products of Free Groups and 3-Manifold Groups, I}
\author{Olga Kharlampovich, Alina Vdovina}
\date{Source: arXiv:2606.00335v3 (CC BY 4.0)}
\begin{document}
\maketitle

\noindent\emph{Source and licence.} Equations in Products of Free Groups and 3-Manifold Groups, I. Version arXiv:2606.00335v3 (CC BY 4.0), distributed under the Creative Commons Attribution 4.0 International licence, as recorded in the arXiv licence metadata for this exact version and shown on the abstract page https://arxiv.org/abs/2606.00335v3. Full rights notice: licenses/arxiv-2606.00335-CC-BY-4.0.txt.\par\smallskip

\noindent\textit{Editorial scope.} Counted unit: the Prop whose source label is \texttt{None} in the article (source0.tex lines 3979--4050, source label \texttt{None}), delivered together with the article's own complete original proof of it (source0.tex lines 4051--4388, one \texttt{proof} environment of 338 lines). No mathematics is written, repaired or re-derived: statement and proof are the article's own text line for line. Required context retained uncounted: none. Every definition, display and label of those lines is retained unchanged. Omitted from this package and not redistributed: 1--3978, 4389--7109. Nothing inside a retained excerpt is omitted or altered: every retained body is delivered exactly as the article's lines are written, so no omission receipt of any kind accompanies this package. Citations: the retained excerpts carry no citation command at all, so no bibliography entry of the article is transcribed and no reference is redistributed. Reference adaptation: none was needed and none was made; every reference inside the delivered unit resolves inside it, so no unresolved reference or citation is produced. Printed slips kept literal: no printed word, symbol, hypothesis or number of the counted statement or of its proof was altered. Layout and class: the article's own class is \texttt{article} with packages including amsmath, amsthm, amsfonts, amssymb, graphicx, ulem, listings, booktabs, longtable, fontenc, beramono, lmodern, color, xcolor; the wrapper is the lane's pinned \texttt{amsart} editorial wrapper at 10pt on A4, which restates the theorem-like environments the retained text needs and adds the article's own macro definitions that the retained text needs (none), with extra wrapper packages (bm, bbm, dsfont, xspace, cancel, mathtools). The delivered numbering is the unit's own. Assets: no retained span contains a figure environment, a graphics-inclusion command, a PDF-page inclusion, a table or a TikZ picture; no image file is copied into this package, the package declares no asset dependency and no image file is redistributed with it. Licence: version arXiv:2606.00335v3 of this article is distributed under the Creative Commons Attribution 4.0 International licence, as recorded in the arXiv licence metadata for this exact version and shown on the abstract page https://arxiv.org/abs/2606.00335v3. A full rights notice is delivered at licenses/arxiv-2606.00335-CC-BY-4.0.txt. Characters: every retained excerpt is pure ASCII, so the delivered engine, XeLaTeX with its default Latin Modern text font, reports no missing character.
\begin{prop}[Two-generator quotient of the genus-four presentation] Let \(K\)  be genus 4 group with exponent sums $(0,0,0)$, exponents $(l,-l),(1,-1),(k-k)$.
\[
\begin{aligned}
K=\big\langle\,&
a_1,a_2,a_3,a_4,\,
x_1,x_2,x_3,x_4,\,
y_1,y_2,y_3,y_4,\\
& z_1,z_2,z_3,z_4,\,
t_1,t_2,t_3,t_4,\,
s_1,s_2,s_3,s_4
\ \big|\ \\[2mm]
& z_1=y_1,\quad
  t_1=z_1,\quad
  t_2=z_2,\\
& y_1=a_1(a_1a_2)^l,\quad
  y_2=a_2(a_1a_2)^l,\quad
  y_3=a_3,\quad
  y_4=a_4,\\
& z_2=y_2(a_2^{y_2}a_3),\quad
  z_3=y_3(a_2^{y_2}a_3),\quad
  z_4=a_4,\\
& t_3=z_3(a_3^{z_3}a_4^{z_4})^k,\quad
  t_4=z_4(a_3^{z_3}a_4^{z_4})^k,\\
& x_1=t_1(a_1^{t_1}a_2^{t_2})^{-l},\quad
  x_2=t_2(a_1^{t_1}a_2^{t_2})^{-l},\\
& x_3=t_3,\quad
  x_4=t_4,\\
& s_1=x_1,\quad
  s_2=x_2(a_2^{x_2}a_3^{x_3})^{-1},\\
& s_3=x_3(a_2^{x_2}a_3^{x_3})^{-1},\quad
  s_4=x_4,\\
& s_1=1,\quad
  s_2=1,\\
& s_3(a_3^{s_3}a_4^{s_4})^{-k}=1,\quad
  s_4(a_3^{s_3}a_4^{s_4})^{-k}=1
\big\rangle .
\end{aligned}
\]

Let
\[
Y=a_1a_2,\qquad
X=a_1^{t_1}a_2^{t_2},\qquad
A=a_2^{y_2}a_3,\qquad
B=a_2^{x_2}a_3^{x_3},
\]
and
\[
C=a_3^{z_3}a_4^{z_4},
\qquad
D=a_3^{s_3}a_4^{s_4}.
\]
 Let
\[
\overline{K}
=
K/\left\langle\!\left\langle X^l,\ D^k\right\rangle\!\right\rangle .
\]
Then \(\overline{K}\) admits the two-generator, four-relator presentation
\[
{
\overline{K}
\cong
\left\langle A,Y \ \middle|\ 
Y^{2l+1},\;
\left(Y^{l+1}A^{-1}Y^{l+1}A\right)^l,\;
\left(A^{-1}Y^lA^2\right)^{2k+1},\;
\left(Y^{-1}A^2\right)^k
\right\rangle .
}\]

\end{prop}

\begin{proof}

First we pass from the original presentation to the six Dehn-twist generators
\[
Y=a_1a_2,\qquad
X=a_1^{t_1}a_2^{t_2},\qquad
A=a_2^{y_2}a_3,\qquad
B=a_2^{x_2}a_3^{x_3},
\]
and
\[
C=a_3^{z_3}a_4^{z_4},
\qquad
D=a_3^{s_3}a_4^{s_4}.
\]

From the relations \(s_1=1\) and \(s_2=1\), we get
\[
x_1=1,\qquad x_2=B.
\]
Since
\[
x_1=t_1X^{-l},\qquad x_2=t_2X^{-l},
\]
it follows that
\[
t_1=X^l,\qquad t_2=BX^l.
\]
But \(t_1=z_1=y_1\), and
\[
y_1=a_1Y^l.
\]
Therefore
\[
a_1Y^l=X^l,
\]
so
\[
a_1=X^lY^{-l}.
\]
Similarly,
\[
t_2=z_2=y_2A.
\]
Since \(y_2=a_2Y^l\), we get
\[
a_2Y^lA=BX^l,
\]
and hence
\[
a_2=BX^lA^{-1}Y^{-l}.
\]

Next, from
\[
s_3D^{-k}=1,\qquad s_4D^{-k}=1,
\]
we get
\[
s_3=D^k,\qquad s_4=D^k.
\]
Since
\[
s_3=x_3B^{-1}=t_3B^{-1}
\]
and
\[
t_3=z_3C^k,
\]
we have
\[
z_3C^kB^{-1}=D^k.
\]
Thus
\[
z_3=D^kBC^{-k}.
\]
But
\[
z_3=y_3A=a_3A,
\]
so
\[
a_3=D^kBC^{-k}A^{-1}.
\]
Likewise,
\[
s_4=x_4=t_4=z_4C^k.
\]
Since \(z_4=a_4\), we get
\[
a_4C^k=D^k,
\]
and hence
\[
a_4=D^kC^{-k}.
\]

Therefore the original presentation is Tietze-equivalent to the
six-generator presentation
\[
\begin{aligned}
K\cong
\langle X,Y,A,B,C,D \mid\;&
Y=X^lY^{-l}BX^lA^{-1}Y^{-l},\\
&
A=Y^{-l}BX^lA^{-1}D^kBC^{-k}A^{-1},\\
&
C=A^{-1}D^kBC^{-k}D^kC^{-k},\\
&
X=Y^{-l}X^lA^{-1}Y^{-l}BX^l,\\
&
B=X^lA^{-1}Y^{-l}BC^{-k}A^{-1}D^kB,\\
&
D=BC^{-k}A^{-1}D^kC^{-k}D^k
\rangle .
\end{aligned}
\]

Now pass to the quotient
\[
\overline K
=
K/\left\langle\!\left\langle X^l,\ D^k\right\rangle\!\right\rangle .
\]
Thus in \(\overline K\) we impose
\[
X^l=1,\qquad D^k=1.
\]
The six relations reduce to
\[
\begin{aligned}
Y&=Y^{-l}BA^{-1}Y^{-l},&\qquad
A&=Y^{-l}BA^{-1}BC^{-k}A^{-1},\\
C&=A^{-1}BC^{-2k},&
X&=Y^{-l}A^{-1}Y^{-l}B,\\
B&=A^{-1}Y^{-l}BC^{-k}A^{-1}B,&
\text{and}\quad D&=BC^{-k}A^{-1}C^{-k}.
\end{aligned}
\]

From the first relation,
\[
Y=Y^{-l}BA^{-1}Y^{-l},
\]
we obtain
\[
Y^{l+1}=BA^{-1}Y^{-l}.
\]
Multiplying on the right by \(Y^l\), we get
\[
Y^{2l+1}=BA^{-1}.
\]
Hence
\[
B=Y^{2l+1}A.
\]

Substituting this into the second relation gives
\[
A
=
Y^{-l}(Y^{2l+1}A)A^{-1}
(Y^{2l+1}A)C^{-k}A^{-1}.
\]
Thus
\[
A=Y^{3l+2}AC^{-k}A^{-1}.
\]
Multiplying on the right by \(A\), and then multiplying on the left by
\((Y^{3l+2}A)^{-1}=A^{-1}Y^{-(3l+2)}\), gives
\[
C^{-k}=A^{-1}Y^{-(3l+2)}A^2.
\]

Now substitute
\[
B=Y^{2l+1}A
\]
and
\[
C^{-k}=A^{-1}Y^{-(3l+2)}A^2
\]
into the fifth relation
\[
B=A^{-1}Y^{-l}BC^{-k}A^{-1}B.
\]
We get
\[
Y^{2l+1}A
=
A^{-1}Y^{-l}(Y^{2l+1}A)
\left(A^{-1}Y^{-(3l+2)}A^2\right)
A^{-1}(Y^{2l+1}A).
\]
The right-hand side simplifies as follows:
\[
\begin{aligned}
& A^{-1}Y^{-l}(Y^{2l+1}A)
\left(A^{-1}Y^{-(3l+2)}A^2\right)
A^{-1}(Y^{2l+1}A) \\
&\qquad =
A^{-1}Y^{l+1}Y^{-(3l+2)}AY^{2l+1}A \\
&\qquad =
A^{-1}Y^{-(2l+1)}AY^{2l+1}A .
\end{aligned}
\]
Therefore
\[
Y^{2l+1}A
=
A^{-1}Y^{-(2l+1)}AY^{2l+1}A.
\]
Canceling the final \(A\) gives
\[
Y^{2l+1}
=
A^{-1}Y^{-(2l+1)}AY^{2l+1}.
\]
Multiplying on the right by \(Y^{-(2l+1)}\), we obtain
\[
1=A^{-1}Y^{-(2l+1)}A.
\]
Hence
\[
Y^{2l+1}=1.
\]
Consequently,
\[
B=A.
\]

Since \(Y^{2l+1}=1\), we have
\[
Y^{-(3l+2)}=Y^l.
\]
Therefore
\[
C^{-k}=A^{-1}Y^lA^2.
\]

The third relation is
\[
C=A^{-1}BC^{-2k}.
\]
Since \(B=A\), this becomes
\[
C=C^{-2k}.
\]
Using
\[
C^{-k}=A^{-1}Y^lA^2,
\]
we get
\[
C=\left(A^{-1}Y^lA^2\right)^2.
\]
Substituting this back into \(C^{-k}=A^{-1}Y^lA^2\), we obtain
\[
\left(\left(A^{-1}Y^lA^2\right)^2\right)^{-k}
=
A^{-1}Y^lA^2.
\]
Equivalently,
\[
\left(A^{-1}Y^lA^2\right)^{-2k}
=
A^{-1}Y^lA^2.
\]
Hence
\[
\left(A^{-1}Y^lA^2\right)^{2k+1}=1.
\]

The fourth relation gives
\[
X=Y^{-l}A^{-1}Y^{-l}B.
\]
Using \(B=A\) and \(Y^{2l+1}=1\), we have
\[
Y^{-l}=Y^{l+1}.
\]
Therefore
\[
X=Y^{l+1}A^{-1}Y^{l+1}A.
\]
Since \(X^l=1\) in \(\overline K\), we obtain the relator
\[
\left(Y^{l+1}A^{-1}Y^{l+1}A\right)^l=1.
\]

Finally, the sixth relation gives
\[
D=BC^{-k}A^{-1}C^{-k}.
\]
Using \(B=A\) and \(C^{-k}=A^{-1}Y^lA^2\), we get
\[
D
=
A(A^{-1}Y^lA^2)A^{-1}(A^{-1}Y^lA^2).
\]
Thus
\[
D=Y^{2l}A^2.
\]
Since \(Y^{2l+1}=1\), this becomes
\[
D=Y^{-1}A^2.
\]
Because \(D^k=1\) in \(\overline K\), we obtain
\[
\left(Y^{-1}A^2\right)^k=1.
\]

We have eliminated \(B,C,X,D\), and the remaining generators are only
\(A\) and \(Y\). The remaining relators are
\[
\begin{aligned}
Y^{2l+1}&=1,&
\left(Y^{l+1}A^{-1}Y^{l+1}A\right)^l&=1,\\
\left(A^{-1}Y^lA^2\right)^{2k+1}&=1,&
\text{and}\quad \left(Y^{-1}A^2\right)^k&=1.
\end{aligned}
\]
Therefore
\[
\overline K
\cong
\left\langle A,Y \ \middle|\ 
Y^{2l+1},\;
\left(Y^{l+1}A^{-1}Y^{l+1}A\right)^l,\;
\left(A^{-1}Y^lA^2\right)^{2k+1},\;
\left(Y^{-1}A^2\right)^k
\right\rangle .
\]
This proves the proposition. 

\end{proof}

\end{document}
